%%% trecho-1
Consulta

%%% trecho-2
Definições e exemplos na notação usada nas contas e nas simulações.

%%% trecho-3
Pesquisar um termo ou exemplo

%%% trecho-5
Ação

%%% trecho-6
Escolha disponível ao agente em um estado. No carrinho, pode-se aplicar força à esquerda ou à direita.

%%% trecho-7
Adam

%%% trecho-8
Otimizador que usa médias móveis do gradiente e de seu quadrado para ajustar a escala das atualizações. Lapan o emprega para treinar a rede do CartPole.

%%% trecho-9
Agente e ambiente

%%% trecho-10
O agente escolhe a ação. O ambiente devolve a próxima observação e a recompensa. No lago, o sorteio do deslizamento pertence ao ambiente.

%%% trecho-11
Corte e percentil

%%% trecho-12
Um percentil localiza uma posição na amostra ordenada. O corte inclusivo aceita empates no limiar. Se o percentil 70 for zero, todos os retornos não negativos passam.

%%% trecho-13
Desconto ($\gamma$)

%%% trecho-14
Peso relativo de um passo de espera. Se $\gamma=0{,}9$, uma recompensa unitária recebida após três ações contribui 0,81 para $G_0$.

%%% trecho-15
Elite

%%% trecho-16
Episódios aceitos pelo corte do lote. Todas as decisões desses episódios entram no ajuste, inclusive uma ação ocasionalmente ruim.

%%% trecho-17
Entropia cruzada

%%% trecho-18
Média do custo $-\ln\pi(a\mid s)$ sob a distribuição de referência. Atribuir probabilidade 0,1 à ação registrada custa aproximadamente 2,30.

%%% trecho-19
Episódio e trajetória

%%% trecho-20
O episódio é uma tentativa completa. A trajetória registra a sequência de estados, ações e recompensas. Chegar a um buraco encerra um episódio do lago.

%%% trecho-21
Esperança

%%% trecho-22
Média teórica que pondera cada resultado pela sua probabilidade. Se os únicos retornos são 0 e 1, o retorno esperado é a probabilidade de sucesso.

%%% trecho-23
Estado e observação

%%% trecho-24
O estado reúne a informação que governa a evolução do ambiente. A observação é a informação entregue ao agente. No CartPole deste trabalho, observam-se as quatro coordenadas do estado.

%%% trecho-25
Gradiente

%%% trecho-26
Vetor das derivadas parciais da perda. Na saída softmax, a derivada por logit é a probabilidade prevista menos o indicador da ação observada.

%%% trecho-27
Logit e softmax

%%% trecho-28
Logit é a pontuação de uma ação antes da normalização. A softmax exponencia as pontuações e divide pela soma. Os logits $(0,\ln 2)$ dão probabilidades $(1/3,2/3)$.

%%% trecho-29
Lote

%%% trecho-30
Conjunto de episódios coletados antes de uma atualização. Nas medições em Python, cada lote contém dezesseis tentativas.

%%% trecho-31
Política ($\pi$)

%%% trecho-32
Distribuição das ações dado o estado. Uma política uniforme no lago atribui probabilidade $1/4$ a cada uma das quatro ações.

%%% trecho-33
Recompensa ($R$)

%%% trecho-34
Número devolvido após uma ação. Em FrozenLake, a chegada à meta vale 1 e uma transição comum vale 0.

%%% trecho-35
Retorno ($G$)

%%% trecho-36
Soma das recompensas futuras, com o desconto escolhido. Com recompensa de um ponto por passo e $\gamma=1$, um episódio de 28 passos vale 28.

%%% trecho-37
Semente

%%% trecho-38
Inteiro que inicializa o gerador pseudoaleatório. Repetir a configuração e a semente no mesmo simulador permite repetir a sequência de sorteios.

%%% trecho-39
Verossimilhança

%%% trecho-40
Probabilidade atribuída aos dados observados, considerada como função dos parâmetros. Para D, E, D em um estado, vale $p^2(1-p)$.

%%% termo-0
Amostragem por importância

%%% def-0
Estima uma esperança sob $P$ usando amostras de $q$ e pesos $P/q$. A distribuição de coleta precisa cobrir os resultados relevantes.

%%% termo-1
Decisões aceitas ($M$)

%%% def-1
Total de passos dos episódios de elite. Dois episódios com 6 e 10 passos dão $M=16$ rótulos.

%%% termo-2
Desempenho esperado ($J(\theta)$)

%%% def-2
Média teórica dos retornos em trajetórias geradas com os pesos atuais.

%%% termo-3
Divergência de Kullback--Leibler (KL)

%%% def-3
Para distribuições discretas, $D_{\mathrm{KL}}(p\Vert q)=\sum_x p(x)\ln[p(x)/q(x)]$. Vale zero quando coincidem e pode depender da ordem.

%%% termo-4
Fração nominal de elite ($\rho$)

%%% def-4
Proporção usada para definir o corte. Com $\rho=0{,}3$, usa-se o percentil 70. Empates podem alterar a proporção aceita.

%%% termo-5
Indicador

%%% def-5
Função que vale um quando a condição é satisfeita e zero caso contrário. Por exemplo, $\mathbf 1\{A_t=b\}$.

%%% termo-6
Limiar ($c_k$)

%%% def-6
Valor do corte calculado a partir dos escores do lote de índice $k$.

%%% termo-7
On-policy

%%% def-7
Coleta feita com a política atual. A versão básica do CEM joga um novo lote após cada ajuste.

%%% termo-8
Parâmetros ($\theta$)

%%% def-8
Pesos que determinam as preferências por ação. No lago, há uma linha de pesos por casa.

%%% termo-9
Quantil

%%% def-9
Valor que delimita uma fração da amostra ordenada. O quantil linear de ordem 0,5 entre 2 e 6 vale 4.

%%% termo-10
Sem modelo

%%% def-10
O agente aprende das transições observadas sem precisar da distribuição de transição do ambiente.

%%% termo-11
Taxa de aprendizagem ($\alpha$)

%%% def-11
Fator que multiplica o gradiente na atualização dos pesos. Duplicar a taxa duplica o primeiro deslocamento com o mesmo gradiente.
